\documentclass[a4paper]{article} %

\date{} %

% Please, do not change any of the above lines

\usepackage{amssymb} % add other LaTeX packages you need here.

\begin{document}
\setcounter{page}{1} % Please, do not delete this line

% Your title
\title{Mal'tsev conditions arising from functional relations}

% Authors
\author{Gerg\H{o} Gyenizse, Mikl\'os Mar\'oti, L\'aszl\'o Z\'adori \\ %
       University of Szeged\\ \\ % Affiliation 1
       \texttt{gyenizse.gergo@math.u-szeged.hu} % Only one corresponding e-mail, please
       }%

\maketitle

% The abstract

We call a relation $\mathbf R\subseteq A^k$ {\em functional} if it is the graph of an operation i.e. if its first $k-1$ coordinates uniquely determine its last. If $\mathcal V$ is a variety, the {\em free functional relation} generated by $\mathbf R$ in $\mathcal V$ is obtained by factoring out the free relational structure generated by $\mathbf R$ with the smallest congruence $\theta$ such that $\mathbf R\slash\theta$ is functional.

The free functional relation can be trivial. For a given $\mathbf R$, the triviality of the free functional relation generated by $\mathbf R$ in $\mathcal V$ can be naturally described by a Mal'tsev condition. In this talk, we explore some of these Mal'tsev conditions, and the Galois connection between relations and varieties induced by the triviality of the free functional relation. 



\end{document}
















